\section{Оформление}
\subsection{Перечисления}

Ненумерованное перечисление
\begin{itemize}
\item первый элемент;
\item второй элемент;
\item третий элемент.
\end{itemize}

Нумерованное перечисление:
\begin{enumerate}
\item первый элемент;
\item второй элемент состоит из:
\begin{enumerate}
\item первого подэлемента;
\item второго подэлемента;
\end{enumerate}
\item третий элемент.
\end{enumerate}

\subsection{Таблицы}

Название таблицы следует помещать над таблицей. Проверка
таблицы~\ref{t:1}.

\begin{table}[b]
\caption{Заголовок таблицы}
\label{t:1}
\begin{tabular}{|c|c|c|c|}
\hline
\multicolumn{1}{|p{4cm}|}{Диаметр стержня крепежной детали, мм}&
\multicolumn{1}{p{4cm}|}{Масса $1000~\text{шт.}$ стальных шайб, кг}&
\multicolumn{1}{p{4cm}|}{Диаметр стержня крепежной детали, мм}&
\multicolumn{1}{p{4cm}|}{Масса $1000~\text{шт.}$ стальных шайб,
кг}\\\hline
$1{,}1$ & $0{,}045$ & $2{,}0$ & $0{,}192$\\\hline
$1{,}2$ & $0{,}043$ & $2{,}5$ & $0{,}350$\\\hline
$1{,}4$ & $0{,}111$ & $3{,}0$ & $0{,}553$\\\hline
\end{tabular}
\end{table}

\subsection{Иллюстрации}

Пример иллюстрации изображен на рисунке~\ref{pidgin}.

\begin{figure}[ht]
  \center
  \includegraphics{pidgin}
  \caption{Pidgin}\label{pidgin}
\end{figure}

\subsection{Математика}

Плотность каждого образца $\rho, \text{кг}/\text{м}^3$, вычисляют
по формуле
\begin{equation}
\label{eq:1}
\rho = \frac{m}{V},
\end{equation}
\begin{ESKDexplanation}
\item[где ] $m$ "--- масса образца, кг;
\item $V$ "--- объем образца, $\text{м}^3$.
\end{ESKDexplanation}
